\documentclass{article} % For LaTeX2e
\usepackage[preprint]{colm2026_conference}

\usepackage{microtype}
\usepackage{hyperref}
\usepackage{url}
\usepackage{booktabs}
\usepackage{lineno}
\usepackage{amsmath,amsfonts,amssymb}
\usepackage{tikz}
\usetikzlibrary{trees, shapes.geometric, arrows.meta, positioning, calc, shadows, backgrounds, fit}
\usepackage[edges]{forest}
\usepackage{enumitem}
\usepackage{multirow}

\definecolor{darkblue}{rgb}{0, 0, 0.5}
\hypersetup{colorlinks=true, citecolor=darkblue, linkcolor=darkblue, urlcolor=darkblue}

\title{Who Decides the Training Distribution? \\ A Survey of On-Policy Distillation for Large Language Models}

\author{Mingyang Song \& Mao Zheng \\
Large Language Model Department \\Tencent, China\\
\texttt{nickmysong@tencent.com}}

\newcommand{\fix}{\marginpar{FIX}}
\newcommand{\new}{\marginpar{NEW}}

% Math macros
\newcommand{\E}{\mathbb{E}}
\newcommand{\KL}{D_{\mathrm{KL}}}
\newcommand{\ptheta}{p_\theta}
\newcommand{\pteacher}{p_T}
\newcommand{\pdata}{p_{\mathrm{data}}}
\newcommand{\loss}{\mathcal{L}}
\DeclareMathOperator*{\argmin}{arg\,min}
\DeclareMathOperator*{\argmax}{arg\,max}

\begin{document}

\ifcolmsubmission
\linenumbers
\fi

\maketitle


\begin{abstract}
Knowledge distillation (KD) has long served as the primary method for model compression. However, Large Language Models (LLMs) have driven a shift from static, off-policy distillation to On-Policy Distillation (OPD). This survey provides a theoretically grounded overview of OPD, organized around a core insight: the essence of OPD is not the choice of divergence metric, but a shift in \textit{"who decides the training distribution."}  By letting the student's own policy dictate the state distribution, OPD directly mitigates the train-test mismatch and compounding exposure bias of autoregressive decoding.
We taxonomize the field into three trajectories based on feedback signals---Logit-based, Outcome-based, and Self-play---and analyze neglected dimensions including the compute-quality tradeoff, the teacher's ``echo chamber'' in uncertain regions, and dynamic divergence adaptation. Finally, we chart future directions from distillation scaling laws to agent-level OPD.
\end{abstract}

\section{Introduction}
\label{sec:intro}

Large Language Models (LLMs) achieve strong performance across reasoning and generation tasks, but frontier models' massive parameter counts impose constraints on deployment, latency, and energy efficiency. DeepSeek-R1 \citep{2501.12948} demonstrated that the reasoning capabilities of a 671B Mixture-of-Experts model can be distilled into dense models as small as 7B parameters. This result confirms that knowledge distillation is no longer a marginal optimization---it is the central engine for democratizing frontier intelligence on consumer-grade hardware.

Traditionally, distillation in NLP has been dominated by off-policy methods, where a student mimics a teacher's token-level probabilities over a fixed dataset drawn from $\pdata$. As shown by \citet{2305.15717}, off-policy distillation suffers from a fundamental flaw: \textit{train-test mismatch}. During training, the student conditions on gold-standard prefixes (teacher-forcing); during inference, it must condition on its own previously generated tokens. When the student deviates from the training distribution, it enters out-of-distribution territory with no supervisory signal, triggering a cascade of compounding errors known as \textit{exposure bias}. 

To bridge this gap, the field is migrating towards \textbf{On-Policy Distillation (OPD)}. Drawing on Imitation Learning, specifically the DAgger algorithm \citep{ross2011reduction}, OPD has the student generate trajectories under its own policy ($\pi_\theta$), then queries the teacher for feedback on these student-generated paths.
\textbf{Core Insight.} \textit{The essence of On-Policy Distillation is not the choice of divergence (Forward KL vs.\ Reverse KL), but a shift in ``who decides the training distribution.''} The transition from $y \sim \pdata$ to $y \sim \ptheta$ transforms the learning objective from memorizing an alien distribution to navigating and correcting the student's own reachable generative space.

Despite the rapid growth of OPD methods---from GKD \citep{2306.13649} to MiniLLM \citep{2306.08543} and SPIN \citep{2401.01335}---the literature remains fragmented. Existing surveys such as \citet{2402.13116} primarily treat distillation as monolithic model compression, lacking a unified mathematical perspective on the on-policy dynamics governing modern LLM alignment and reasoning transfer. This paper addresses this gap.

Specifically, the contributions of this review are:
\begin{itemize}
    \item \textbf{A Unified Theoretical Framework.} We formulate the transition from off-policy to on-policy distillation via sequential decision-making and f-divergence minimization. This framework unifies methods like GKD \citep{2306.13649}, MiniLLM \citep{2306.08543}, and DistiLLM \citep{2402.03898} as optimizing the same underlying objective under different divergence constraints and sampling mixtures.
    \item \textbf{A Feedback-Centric Taxonomy.} We categorize OPD into three routes based on the supervisory signal: \textit{Logit-based} (dense, white-box), \textit{Outcome-based} (sparse, flexible, using PRMs \citep{2305.20050}), and \textit{Self-play} (teacher-free, prone to saturation). We highlight how these tracks are converging into hybrid methods such as AlignDistil \citep{2503.02832} and MiMo-V2 \citep{2601.02780}.
    \item \textbf{Identification of Neglected Issues.} We surface challenges the community has largely ignored: (a) the \textit{compute-quality tradeoff} where online generation costs may outweigh distillation benefits; (b) the ``echo chamber'' effect from querying teachers in high-uncertainty regions; (c) the need for \textit{dynamic divergence adaptation}; and (d) reconceptualizing distillation as structured knowledge transfer.
    \item \textbf{Bridging White-box and Black-box.} We analyze how on-policy techniques adapt when teacher internals are inaccessible, comparing logit matching with API-driven approaches like Lion \citep{2305.12870} and GAD \citep{2511.10643}.
    \item \textbf{Future Directions.} We chart open problems including distillation scaling laws, uncertainty-aware OPD, and agent-level environmental distillation.
\end{itemize}

The remainder of this survey is organized as follows. Section~\ref{sec:background} establishes the mathematical preliminaries. Section~\ref{sec:taxonomy} presents our taxonomy of OPD methodologies. Subsequent sections dissect logit-based optimization, outcome-reward architectures, reasoning distillation, industrial systems, and open challenges.


\section{Background and Preliminaries}
\label{sec:background}

This section formalizes the sequential nature of LLM generation, defines distribution matching and exposure bias, and constructs a unified framework via the f-divergence family.

\subsection{Knowledge Distillation Fundamentals}
\label{subsec:kd_fundamentals}

Classical Knowledge Distillation \citep{hinton2015distilling} transfers the ``dark knowledge'' of a teacher network $\mathcal{T}$ to a student $\mathcal{S}$. Given logits $z \in \mathbb{R}^{|V|}$, a temperature $\tau > 0$ produces softened probabilities:
\begin{equation}
    P(y|x; \tau) = \frac{\exp(z_y / \tau)}{\sum_{y' \in V} \exp(z_{y'} / \tau)}
\end{equation}
The KD objective minimizes the KL divergence between teacher and student softened distributions:
\begin{equation}
    \loss_{\mathrm{KD}} = \tau^2 \KL(\pteacher(\cdot|x; \tau) \parallel \ptheta(\cdot|x; \tau)) = \tau^2 \sum_{y \in V} \pteacher(y|x; \tau) \log \left( \frac{\pteacher(y|x; \tau)}{\ptheta(y|x; \tau)} \right)
\end{equation}
The $\tau^2$ factor compensates for the $1/\tau^2$ gradient scaling as temperature increases. The gradient with respect to student logits $z_i^{\mathcal{S}}$ simplifies to $\tau (p_i^{\mathcal{S}} - p_i^{\mathcal{T}})$. In the high-temperature limit ($\tau \to \infty$), applying $\exp(z_i/\tau) \approx 1 + z_i/\tau$ shows this reduces to MSE between raw logits: $\frac{1}{|V|} (z_i^{\mathcal{S}} - z_i^{\mathcal{T}})$. Thus high-temperature KD degenerates into logit matching, forcing the student to replicate the teacher's full logit landscape including sub-optimal probability masses (dark knowledge). However, when applied to autoregressive LLMs, this formulation assumes the prefix is drawn from the data distribution, ignoring the sequential, compounding nature of generation.

\subsection{Token-Level vs. Sequence-Level Distillation}
\label{subsec:token_vs_sequence}

Applying \citet{hinton2015distilling}'s formulation to language modeling yields token-level distillation. For a sequence $y = (y_1, \dots, y_T)$ with prefix $y_{<t}$:
\begin{equation}
    \loss_{\mathrm{Token-KD}} = \mathbb{E}_{x, y \sim \mathcal{D}} \left[ \sum_{t=1}^{|y|} \KL(\pteacher(\cdot | x, y_{<t}) \parallel \ptheta(\cdot | x, y_{<t})) \right]
\end{equation}
This is computationally tractable but misaligned with the global generative objective: it optimizes next-token accuracy assuming the history $y_{<t}$ is error-free.

\citet{kim2016sequence} proposed Sequence-Level KD to match joint distributions over the full sequence space:
\begin{equation}
    \loss_{\mathrm{Seq-KD}} = \KL(P_{\mathcal{T}}(y | x) \parallel P_{\theta}(y | x)) = \sum_{y \in \mathcal{Y}} P_{\mathcal{T}}(y | x) \log \left( \frac{P_{\mathcal{T}}(y | x)}{P_{\theta}(y | x)} \right)
\end{equation}
Since $|\mathcal{Y}| = |V|^T$ makes exact computation intractable, they approximate $P_{\mathcal{T}}(y|x)$ with a Dirac delta at the teacher's beam-search output $\hat{y}$, collapsing the loss to:
\begin{equation}
    \loss_{\mathrm{Seq-Approx}} = - \sum_{t=1}^{|\hat{y}|} \log \ptheta(\hat{y}_t | x, \hat{y}_{<t})
\end{equation}
This forces the student to model complete teacher-favored trajectories, implicitly transferring sequence-level dependencies. However, Sequence-Level KD remains \textit{off-policy}: the student still trains via teacher-forcing on static trajectories and never conditions on its own prefixes. It therefore remains susceptible to the train-test mismatch it sought to resolve.

\subsection{The Distribution Mismatch Problem and Exposure Bias}
\label{subsec:exposure_bias}

The pathology motivating OPD is \textit{exposure bias}, a covariate shift in sequential decision-making. We cast autoregressive generation as an MDP: states $s_t = (x, y_{<t})$ are token prefixes, actions are vocabulary tokens, and transitions are deterministic. During off-policy training, states follow the data distribution $d_{\pdata}(s)$:
\begin{equation}
    \loss_{\mathrm{train}} = \mathbb{E}_{s \sim d_{\pdata}} \left[ \KL(\pteacher(\cdot|s) \parallel \pi_\theta(\cdot|s)) \right]
\end{equation}
At inference, the student follows its own state distribution $d_{\ptheta}(s) \neq d_{\pdata}(s)$. The student receives zero signal in states within $d_{\ptheta}(s) \setminus d_{\pdata}(s)$.

The severity is quantified by the DAgger analysis \citep{ross2011reduction}: if per-step error under the training distribution is $\epsilon$, the total discrepancy over length $T$ under the student's distribution scales as:
\begin{equation}
    \mathbb{E}_{s \sim d_{\ptheta}} \left[ \sum_{t=1}^T \loss(s_t) \right] \le O(\epsilon T^2)
\end{equation}
This $O(T^2)$ bound is severe for LLMs generating thousands of tokens: a single error shifts the prefix out-of-distribution, increasing the probability of subsequent errors. OPD resolves this by changing the training expectation from $\mathbb{E}_{s \sim d_{\pdata}}$ to $\mathbb{E}_{s \sim d_{\ptheta}}$, forcing the student to confront its own mistakes and learn recovery behaviors, reducing error accumulation to $O(\epsilon T)$.

\subsection{The f-Divergence Family}
\label{subsec:f_divergence}

OPD methods are parameterized by the choice of f-divergence. For distributions $P$ and $Q$ and convex $f$ with $f(1)=0$:
\begin{equation}
    D_f(P \parallel Q) = \mathbb{E}_{x \sim Q} \left[ f\left(\frac{P(x)}{Q(x)}\right) \right]
\end{equation}
The key members relevant to LLM distillation are:

\textbf{Forward KL} ($f(u) = u \log u$): $D_{\mathrm{KL}}(P \parallel Q) = \mathbb{E}_{x \sim P}[\log(P(x)/Q(x))]$. The student must cover all teacher modes (\textit{mode-covering}/zero-avoiding). When the student lacks capacity to cover the teacher's full distribution, it places mass between valid modes, producing average or nonsensical outputs.

\textbf{Reverse KL} ($f(u) = -\log u$): $D_{\mathrm{KL}}(Q \parallel P) = \mathbb{E}_{x \sim Q}[\log(Q(x)/P(x))]$. The student only places mass where the teacher supports it (\textit{mode-seeking}/zero-forcing). It collapses onto one high-probability mode, which is desirable for generation: producing one coherent answer is preferable to an amalgamation of several.

\textbf{Jensen-Shannon Divergence}: Measures divergence to the mixture $M = \frac{1}{2}(P+Q)$, bounded in $[0, \log 2]$, balancing mode-seeking and mode-covering with stable gradients.

\textbf{Total Variation Distance} ($f(u) = \frac{1}{2}|u-1|$): $\sup_A |P(A) - Q(A)|$. Provides robust distance measurement regardless of extreme log-probability ratios, though non-differentiability at $u=1$ complicates optimization.

\subsection{A Unified View of Modern OPD}
\label{subsec:unified_view}

Recent OPD methods can be unified under a single objective by decoupling the sampling distribution from the divergence metric:
\begin{equation}
    \loss_{\mathrm{OPD}}(\theta) = \mathbb{E}_{y \sim \pi_{\mathrm{mix}}} \left[ \sum_{t=1}^{|y|} \mathcal{D}_f \left( \pteacher(\cdot | x, y_{<t}) \parallel \ptheta(\cdot | x, y_{<t}) \right) \right]
\end{equation}
where $\pi_{\mathrm{mix}}$ is the behavioral policy and $\mathcal{D}_f$ is a chosen f-divergence. The major methods map as follows:

\textbf{GKD \citep{2306.13649}:} Defines $\pi_{\mathrm{mix}}$ as an interpolation---with probability $\lambda$ prefixes come from dataset $\mathcal{D}$, with probability $1-\lambda$ from $\ptheta$. Flexible in $\mathcal{D}_f$ (Forward KL, Reverse KL, JSD). Setting $\lambda \to 0$ yields purely on-policy training.

\textbf{MiniLLM \citep{2306.08543}:} Sets $\mathcal{D}_f = D_{\mathrm{KL}}(\ptheta \parallel \pteacher)$ (Reverse KL) with $\pi_{\mathrm{mix}} = \ptheta$. Since the student appears in both the expectation and the log-ratio, optimization requires REINFORCE:
\begin{equation}
    \nabla_\theta \loss_{\mathrm{MiniLLM}} = \mathbb{E}_{y \sim \ptheta} \left[ \sum_{t=1}^{|y|} \left( \log \frac{\ptheta(y_t|y_{<t})}{\pteacher(y_t|y_{<t})} - 1 \right) \nabla_\theta \log \ptheta(y_t|y_{<t}) \right]
\end{equation}
MiniLLM uses value networks as baselines to reduce the high variance of policy gradients.

\textbf{DistiLLM \citep{2402.03898}:} Maintains $\pi_{\mathrm{mix}} = \ptheta$ but reparameterizes the divergence via a skewed distribution $\tilde{P}_{\mathcal{T}} = \alpha \pteacher + (1-\alpha) \ptheta$. Minimizing Forward KL toward this skewed target mimics Reverse KL geometry while allowing standard cross-entropy optimization without policy gradients.

This unified perspective shows that the evolution of OPD is a systematic tightening of bounds on trajectory sampling and distribution matching, not a series of ad hoc algorithms.


\section{Taxonomy of On-Policy Distillation}
\label{sec:taxonomy}

To navigate the expanding OPD literature, we propose a taxonomy along three dimensions: (1) \textbf{Feedback Signal}---what information flows from teacher to student; (2) \textbf{Teacher Access}---deployment constraints on the teacher; and (3) \textbf{Granularity}---temporal resolution of the distillation loss. The taxonomy is visualized in Figure~\ref{fig:taxonomy_tree}.

\begin{figure*}[t]
\centering
\definecolor{rootblue}{HTML}{2C3E6B}
\definecolor{dimteal}{HTML}{0D7377}
\definecolor{dimviolet}{HTML}{6C4FA0}
\definecolor{dimorange}{HTML}{C2722A}
\tikzset{
  rbox/.style={fill=rootblue, text=white, rounded corners=10pt,
    font=\normalsize\bfseries\sffamily, minimum height=1cm, text centered,
    inner xsep=10pt},
  dimA/.style={fill=dimteal, text=white, rounded corners=8pt,
    font=\small\bfseries\sffamily, minimum height=0.75cm,
    text centered, inner xsep=8pt},
  dimB/.style={fill=dimviolet, text=white, rounded corners=8pt,
    font=\small\bfseries\sffamily, minimum height=0.75cm,
    text centered, inner xsep=8pt},
  dimC/.style={fill=dimorange, text=white, rounded corners=8pt,
    font=\small\bfseries\sffamily, minimum height=0.75cm,
    text centered, inner xsep=8pt},
  leaf/.style={fill=#1!6, draw=#1!25, rounded corners=4pt,
    font=\footnotesize, inner xsep=6pt, inner ysep=5pt, align=left}
}
\resizebox{\textwidth}{!}{
\begin{forest}
  for tree={
    grow=east,
    reversed=true,
    growth parent anchor=east,
    parent anchor=east,
    child anchor=west,
    l sep=12mm,
    s sep=4mm,
    edge={draw=gray!35, line width=0.8pt, -{Stealth[length=4pt, width=3.5pt, round]}},
    edge path={
      \noexpand\path[\forestoption{edge}]
      (!u.parent anchor) -- +(10pt,0) |- (.child anchor)\forestoption{edge label};
    },
    anchor=west
  }
  [\textbf{On-Policy Distillation}, rbox, text width=2.8cm
    [{Feedback Signal}, dimA, edge={draw=dimteal!60, line width=1.2pt}
      [{\textbf{Logit-Based} --- \textit{GKD, DistiLLM, ToDi, AKL, DSKD, Delta KD, Entropy-Aware}}, leaf=dimteal, edge={draw=dimteal!30}]
      [{\textbf{Outcome-Based} --- \textit{RLKD, AlignDistil, DPO, Reward Extrapolation}}, leaf=dimteal, edge={draw=dimteal!30}]
      [{\textbf{Self-Play} --- \textit{SPIN, OPSD, Online Experiential Learning}}, leaf=dimteal, edge={draw=dimteal!30}]
    ]
    [{Teacher Access}, dimB, edge={draw=dimviolet!60, line width=1.2pt}
      [{\textbf{White-Box} --- \textit{GKD, MiniLLM, DistiLLM, ToDi, PACED, FastOPD}}, leaf=dimviolet, edge={draw=dimviolet!30}]
      [{\textbf{Black-Box} --- \textit{GAD, Lion, Zephyr, ORPO-Distill}}, leaf=dimviolet, edge={draw=dimviolet!30}]
      [{\textbf{Self-Distillation} --- \textit{SPIN, OPSD, MTP, Context Distillation}}, leaf=dimviolet, edge={draw=dimviolet!30}]
    ]
    [{Granularity}, dimC, edge={draw=dimorange!60, line width=1.2pt}
      [{\textbf{Token-Level} --- \textit{GKD, DistiLLM, ToDi, AKL, Entropy-Aware}}, leaf=dimorange, edge={draw=dimorange!30}]
      [{\textbf{Sequence-Level} --- \textit{MiniLLM, Constrained OPD, RLKD}}, leaf=dimorange, edge={draw=dimorange!30}]
    ]
  ]
\end{forest}
}
\caption{Taxonomy of on-policy distillation for large language models. The methodology space is organized along three orthogonal dimensions with representative methods. A single method may appear across multiple categories as these axes are independent.}
\label{fig:taxonomy_tree}
\end{figure*}

\subsection{Feedback Signal}
\label{subsec:tax_feedback}

The most critical differentiator among OPD algorithms is the nature of the supervisory signal.

\textbf{Logit-Based Feedback.} The student's generated sequence is evaluated token-by-token against the teacher's probability distribution. Methods like GKD \citep{2306.13649} and MiniLLM \citep{2306.08543} compute f-divergences at every step. More recent work like ToDi \citep{2505.16297} dynamically re-weights token importance based on long-term trajectory impact. Logit-based feedback provides the densest gradient signal ($|V|$-dimensional at each step) but requires evaluating $\pteacher(\cdot|x, \hat{y}_{<t})$.

\textbf{Outcome-Based Feedback.} As models scale, computing full logits becomes prohibitive. Outcome-based methods use scalar rewards evaluating the quality of student trajectories. Frameworks like RLKD \citep{2505.16142} and PRMs \citep{2305.20050} let the student complete a full response, then query a reward model for a holistic score. Optimization proceeds via PPO or DPO: $\max_\theta \mathbb{E}_{y \sim \ptheta} [R(x, y)] - \beta \KL(\ptheta \parallel \pi_{\mathrm{ref}})$. This gives the student freedom to find novel paths without being tethered to the teacher's phrasing.

\textbf{Self-Play Feedback.} Teacher-free methods construct feedback from the student's evolving policy. SPIN \citep{2401.01335} and OPSD \citep{2601.18734} simulate a game where the current model distinguishes its own generations from reference demonstrations. As the student improves, its negative samples become more challenging, providing a natural curriculum.

\subsection{Teacher Access}
\label{subsec:tax_access}

The availability of teacher internals constrains the allowable mathematical formulations.

\textbf{White-Box Distillation.} With full teacher weights accessible, algorithms perform exact divergence matching. DistiLLM \citep{2402.03898} leverages the full logit tensor for skewed KL optimization. White-box access exploits dark knowledge---the small probabilities on incorrect tokens carry structural regularization---but requires maintaining a massive teacher in memory alongside the student.

\textbf{Black-Box Distillation.} Closed-source frontier models (e.g., GPT-4, Claude) restrict access to API-level outputs. Methods like Zephyr \citep{2310.16944} and Lion \citep{2305.12870} use the teacher as a preference annotator: the student generates candidates, the teacher ranks them, and the student optimizes via contrastive losses. GAD \citep{2511.10643} trains a discriminator to distinguish student and teacher text, sidestepping the need for logits.

\textbf{Self-Distillation.} When no superior teacher exists, the model bootstraps from a historical checkpoint or ensembled version of itself. By generating on-policy responses and filtering via self-evaluation (as in SPIN \citep{2401.01335}), the model refines its policy without external dependencies.

\subsection{Granularity}
\label{subsec:tax_granularity}

\textbf{Token-Level.} The loss is computed at each timestep $t$, providing stable gradients and fast convergence but suffering from myopia: a locally improbable token may be optimal for an alternative reasoning path.

\textbf{Sequence-Level.} Supervision is applied at sequence termination, allowing structural freedom but producing sparse gradients. Credit assignment becomes severe: a 1000-token proof failing at step 999 receives a uniform negative signal, ignoring 998 correct steps.

\textbf{Trajectory-Level (Step-by-Step).} State-of-the-art reasoning distillation chunks trajectories into intermediate steps. Frameworks like StepByStep \citep{2305.02301} and SuperCorrect \citep{2410.09008} use PRMs \citep{2305.20050} to parse reasoning steps. The loss is structured as $\loss = \sum_{k=1}^K \log \ptheta(s_k | s_{<k}) \cdot \hat{A}(s_k)$, where $s_k$ is a trajectory chunk and $\hat{A}$ is the advantage. This balances structural freedom with logical constraint, providing the scaffolding needed to distill reasoning from models like DeepSeek-R1 into compact architectures.


\section{White-Box On-Policy Methods}
\label{sec:white_box}

White-box OPD methods require full access to the teacher's probability distributions. Their defining advantage is dense, token-level feedback; their cost is the memory overhead of maintaining the teacher alongside the student during training.

\subsection{Token-Level Divergence Minimization}
\label{subsec:token_level}

Token-level methods let the student generate trajectories $y_{<t} \sim p_S$, then minimize the divergence between student and teacher distributions at each step. The choice of divergence dictates behavior, stability, and convergence.

\paragraph{Generalized Knowledge Distillation (GKD)}
\citet{2306.13649} introduced GKD as a unified on-policy framework bridging knowledge distillation and imitation learning. The objective is:
\begin{equation}
    \mathcal{L}_{GKD} = \mathbb{E}_{x \sim \mathcal{D}, y \sim p_S(\cdot|x)} \left[ \sum_{t=1}^{|y|} \lambda D\big( p_T(\cdot|x, y_{<t}) \parallel p_S(\cdot|x, y_{<t}) \big) + (1-\lambda) \mathcal{L}_{NLL} \right]
    \label{eq:gkd}
\end{equation}
where $D$ can be Forward KL, Reverse KL, or JSD. GKD significantly mitigates compounding errors but incurs substantial computational overhead from online sampling, and a fixed $\lambda$ cannot account for varying token difficulty.

\paragraph{DistiLLM}
\citet{2402.03898} address the high gradient variance from near-zero token probabilities under Reverse KL by introducing Skew KL with adaptive temperature:
\begin{equation}
    \mathcal{L}_{DistiLLM} = \mathbb{E}_{y \sim p_S} \left[ \sum_t D_{KL}\big(p_T(\cdot|y_{<t}) \parallel \alpha p_T(\cdot|y_{<t}) + (1-\alpha) p_S(\cdot|y_{<t})\big) \right]
    \label{eq:distillm}
\end{equation}
The smoothing parameter $\alpha \in (0,1)$ avoids the zero-division problem of Reverse KL when $p_S \to 0$, yielding faster and more stable convergence than JSD. However, the smoothing may mask low-probability knowledge important for complex reasoning.

\paragraph{ToDi}
\citet{2505.16297} observe that different positions in a sequence (e.g., critical reasoning pivots vs.\ syntactic fillers) demand different matching strictness. ToDi dynamically allocates f-divergences per token:
\begin{equation}
    \mathcal{L}_{ToDi} = \mathbb{E}_{y \sim p_S} \left[ \sum_t \omega_t D_{f^{(t)}} \big( p_T(y_t|y_{<t}) \parallel p_S(y_t|y_{<t}) \big) \right]
    \label{eq:todi}
\end{equation}
Reasoning tokens receive strict Reverse KL while stylistic tokens receive relaxed Forward KL. The limitation is the heuristic computation of importance weights $\omega_t$.

\paragraph{Adaptive KL (AKL) and Delta KD}
AKL \citep{2404.02657} dynamically modulates KL penalty strength based on student uncertainty, amplifying teacher guidance in unexplored regions.
Delta KD \citep{2509.14526} distills only the \textit{delta} between a base pre-trained model and the fine-tuned teacher:
\begin{equation}
    \mathcal{L}_{Delta} = D_{KL}\big( (p_T - p_{base}) \parallel (p_S - p_{base}) \big)
\end{equation}
This filters out already-assimilated knowledge, improving efficiency, but requires simultaneous inference of three models (teacher, student, base).

\paragraph{Entropy-Aware OPD}
\citet{2603.07079} identify the ``Echo Chamber'' effect: when the student generates OOD trajectories, teacher entropy $\mathcal{H}(p_T)$ spikes to near-uniform noise, and fitting this noise degrades performance. Their solution gates distillation by teacher certainty:
\begin{equation}
    \mathcal{L}_{Ent-OPD} = \mathbb{E}_{y \sim p_S} \left[ \sum_t \frac{1}{1 + \mathcal{H}(p_T(\cdot|y_{<t}))} D_{KL}(p_T \parallel p_S) \right]
\end{equation}
This severs the propagation of hallucinated knowledge, making it one of the most robust token-level strategies against distribution shift.

\paragraph{Dual Space Knowledge Distillation (DSKD)}
\citet{2504.11426} conduct on-policy distillation simultaneously in logit and hidden-state spaces:
\begin{equation}
    \mathcal{L}_{DSKD} = \lambda_1 \mathcal{L}_{Logit}(p_T, p_S) + \lambda_2 \sum_{l} \text{MSE}\left( W_{proj} h_S^{(l)}, h_T^{(l)} \right)
\end{equation}
By aligning intermediate representations, DSKD ensures the student learns \textit{how} to represent concepts, not just \textit{what} to output, enhancing generalization on reasoning tasks at the cost of additional hyperparameters for layer mapping and projection initialization.

\textbf{Divergence Choice and Task Alignment.} The choice of divergence encodes strong inductive biases. Forward KL is mean-seeking: it forces the student to cover all teacher modes, making it suitable for creative tasks where mode-dropping is harmful. Reverse KL is mode-seeking: it penalizes the student for generating tokens the teacher deems unlikely, making it superior for mathematical reasoning, coding, and safety alignment where suppressing hallucinations is prioritized over diversity.

\subsection{Sequence-Level On-Policy Methods}
\label{subsec:sequence_level}

Token-level methods provide dense supervision but suffer from greedy optimization, ignoring sequence-level joint probability.

\paragraph{MiniLLM and the REINFORCE Derivation}
MiniLLM \citep{2306.08543} optimizes Sequence-level Reverse KL without token-level teacher forcing.
The sequence-level Reverse KL is:
\begin{equation}
    D_{KL}(p_S \parallel p_T) = \mathbb{E}_{y \sim p_S(\cdot|x)} \left[ \log p_S(y|x) - \log p_T(y|x) \right]
\end{equation}
Since $p_S$ is parameterized by $\theta$, the gradient requires the REINFORCE log-derivative trick:
\begin{align}
    \nabla_\theta D_{KL}(p_S \parallel p_T) &= \nabla_\theta \sum_y p_S(y|x) \left( \log p_S(y|x) - \log p_T(y|x) \right) \nonumber \\
    &= \sum_y \nabla_\theta p_S(y|x) \left( \log \frac{p_S(y|x)}{p_T(y|x)} \right) + \sum_y p_S(y|x) \nabla_\theta \left( \log p_S(y|x) - \log p_T(y|x) \right) \nonumber \\
    &= \sum_y p_S(y|x) \nabla_\theta \log p_S(y|x) \left( \log \frac{p_S(y|x)}{p_T(y|x)} \right) + \sum_y p_S(y|x) \nabla_\theta \log p_S(y|x) - 0 \nonumber \\
    &= \mathbb{E}_{y \sim p_S} \left[ \left( \log p_S(y|x) - \log p_T(y|x) \right) \nabla_\theta \log p_S(y|x) \right] + \nabla_\theta \sum_y p_S(y|x)
\end{align}
Since $\sum_y p_S(y|x) = 1$, the last term vanishes:
\begin{equation}
    \nabla_\theta \mathcal{L}_{MiniLLM} = \mathbb{E}_{y \sim p_S} \left[ \Big( \log p_S(y|x) - \log p_T(y|x) \Big) \nabla_\theta \log p_S(y|x) \right]
    \label{eq:minillm_final}
\end{equation}
This shows that Sequence-Level Reverse KL is equivalent to Policy Gradient RL with reward $r(x, y) = \log p_T(y|x) - \log p_S(y|x)$.

\paragraph{Constrained Optimization View}
The REINFORCE estimator's high variance frequently causes policy collapse. \citet{2509.22921} reformulate sequence-level OPD as a trust-region problem, constraining KL divergence per update step (analogous to PPO) so that $\pi_{k+1}$ stays within an $\epsilon$-ball of $\pi_k$.

\paragraph{Token-level vs. Sequence-level Tradeoff}
Token-level methods optimize an upper bound of sequence-level divergence. They are stable and low-variance but myopic, potentially perfecting local syntax while failing at global coherence. Sequence-level methods capture holistic reasoning structure but require many more training iterations and sophisticated baselines due to high Monte Carlo variance.

\subsection{Hybrid and Adaptive Approaches}
\label{subsec:hybrid}

Pure token-level distillation lacks long-term planning; pure sequence-level distillation suffers from unmanageable variance. Hybrid methods address both while tackling the central bottleneck: the compute cost of on-policy generation.

\paragraph{Fast OPD from Reasoning Prefixes}
Fast OPD \citep{2602.15260} injects teacher-generated reasoning prefixes of varying lengths, restricting the student's on-policy exploration to the suffix. This shortens the RL horizon, reducing gradient variance while still allowing exploration of divergent conclusions.

\paragraph{PromptKD}
PromptKD \citep{2402.12842} prepends engineered meta-prompts to the teacher during distillation, shaping the teacher's distribution to be more structurally predictable and narrowing the capacity gap before divergence minimization begins.

\paragraph{PACED}
\citet{2603.11178} observe that on-policy gradients at the extremes of student capability are wasted: easy sequences yield near-zero gradients, while impossible sequences inject noise. PACED identifies the student's ``Competence Frontier''---the boundary where success probability is above chance but not deterministic---and samples exclusively from it, achieving 3$\times$ distillation efficiency without perplexity degradation.

\paragraph{$\mathcal{X}$-KD: Experiential Knowledge Distillation}
$\mathcal{X}$-KD \citep{2602.12674} maintains an episodic memory of historical error trajectories. The student re-generates from the exact token where it previously diverged from the teacher, concentrating exploration in the regions where the student is weakest.

\subsection{Theoretical Analysis and Method Comparison}
\label{subsec:theory_whitebox}

\citet{2404.02657} and \citet{2402.11890} analyze Forward vs.\ Reverse KL under on-policy generation. Forward KL (zero-avoiding, mean-seeking) causes a capacity-limited student to bridge gaps between valid teacher modes, producing hallucinations. Reverse KL (zero-forcing, mode-seeking) concentrates on one mode, structurally preventing hallucinations at the cost of diversity. \citet{2307.15190} generalize this via a unified f-divergence perspective, showing that continuously modulating the convexity of $f$ enables smooth interpolation between mode-seeking and mode-covering behavior.

A comparison of white-box on-policy methods is presented in Table \ref{tab:white_box_comparison}.

\begin{table*}[tbp]
    \centering
    \caption{Comprehensive comparison of white-box on-policy distillation methods.}
    \label{tab:white_box_comparison}
    \resizebox{\textwidth}{!}{
    \begin{tabular}{@{}lclllll@{}}
        \toprule
        \textbf{Method} & \textbf{Year} & \textbf{Granularity} & \textbf{Divergence} & \textbf{Access} & \textbf{On-Policy Formulation} & \textbf{Key Innovation} \\
        \midrule
        GKD \citep{2306.13649}            & 2023 & Token  & F-KL / R-KL / JSD  & White-box & Student Rollouts     & Unified OPD framework \\
        MiniLLM \citep{2306.08543}        & 2023 & Seq.   & Reverse KL          & White-box & REINFORCE            & Sequence-level reward \\
        DistiLLM \citep{2402.03898}       & 2024 & Token  & Skew KL             & White-box & Interpolated         & $\alpha$-smoothing for Reverse KL \\
        PromptKD \citep{2402.12842}       & 2024 & Hybrid & Forward KL          & White-box & Prompted Generation  & Context-bridged capacity gap \\
        AKL \citep{2404.02657}            & 2024 & Token  & Adaptive KL         & White-box & Uncertainty-weighted & Dynamic penalty scaling \\
        DSKD \citep{2504.11426}           & 2025 & Token  & Logit + Hidden      & White-box & Dual-Space Matching  & Geometric feature transfer \\
        Delta KD \citep{2509.14526}       & 2025 & Token  & Delta Logit         & White-box & Base Subtracted      & Fine-tuned signal only \\
        ToDi \citep{2505.16297}           & 2025 & Token  & Dynamic $f$-div     & White-box & Token-Significance   & Per-token divergence selection \\
        Constrained \citep{2509.22921}    & 2025 & Seq.   & Constrained KL      & White-box & Trust-Region         & REINFORCE variance bounding \\
        Entropy-Aware \citep{2603.07079}  & 2026 & Token  & Entropy-Gated KL    & White-box & Confidence-weighted  & Echo chamber defense \\
        FastOPD \citep{2602.15260}        & 2026 & Hybrid & R-KL on Suffix      & White-box & Prefix-Truncated     & Reasoning prefix reuse \\
        PACED \citep{2603.11178}          & 2026 & Hybrid & Adaptive            & White-box & Frontier Curriculum  & Student competence boundary \\
        $\mathcal{X}$-KD \citep{2602.12674} & 2026 & Token & Counterfactual KL  & White-box & Error-Replay Memory  & Historical divergence focus \\
        \bottomrule
    \end{tabular}
    }
\end{table*}

\section{Black-Box and Self-Distillation Methods}
\label{sec:black_box}

Proprietary models do not expose internal weights or logit distributions. OPD must therefore adapt to settings where the teacher is an opaque oracle. As models exhaust human-curated data, self-distillation---where the model refines its own policy---becomes increasingly important.

\subsection{Black-Box On-Policy Distillation}
\label{subsec:black_box}

In black-box OPD, $p_T$ is inaccessible; the teacher serves only as a scorer or preference ranker over student-generated trajectories.

\paragraph{Generative Adversarial Distillation (GAD)}
\citet{2511.10643} map black-box distillation onto a GAN framework: the student generates sequences, and a discriminator distinguishes student text from teacher outputs. The student optimizes via policy gradients to maximize discriminator confusion, entirely bypassing logit matching at the cost of adversarial training instability.

\paragraph{Lion}
Lion \citep{2305.12870} implements a closed-loop system where the student proposes instructions at its uncertainty frontier, generates responses, and queries the teacher on the same instructions. The resulting quality differential creates a dynamic curriculum targeting the gap between student and teacher.

\paragraph{Zephyr and dDPO}
Zephyr \citep{2310.16944} introduced distilled DPO (dDPO). The student generates response pairs $(y_1, y_2) \sim p_S(\cdot|x)$, the teacher judges which is superior, and the student optimizes:
\begin{equation}
    \mathcal{L}_{dDPO} = - \mathbb{E}_{(x, y_w, y_l)} \left[ \log \sigma \left( \beta \log \frac{p_S(y_w|x)}{p_{ref}(y_w|x)} - \beta \log \frac{p_S(y_l|x)}{p_{ref}(y_l|x)} \right) \right]
\end{equation}
This distills the teacher's implicit preferences without logit access or explicit reward modeling.

\paragraph{ORPO-Distill}
\citet{2509.25100} eliminate the need for a frozen reference model by fusing cross-entropy alignment with preference distillation via odds-ratio penalization, reducing the memory footprint of on-policy generation.

\subsection{Self-Play and Self-Distillation}
\label{subsec:self_distill}

When external teachers are unavailable, models must generate their own distillation targets.

\paragraph{SPIN}
\citet{2401.01335} construct a self-play framework where the model at iteration $t$ generates trajectories and the model at $t-1$ serves as baseline. The objective pushes new generations away from the historical policy while pulling toward ground-truth. Convergence analysis shows that with sufficient expressiveness, the SPIN objective's global minimum ensures $p_\theta = p_{data}$.

\paragraph{OPSD}
OPSD \citep{2601.18734} leverages test-time compute for reasoning: the model generates many CoT traces, a rule-based verifier filters correct ones, and the model distills these verified paths back into its parameters, converting inference-time search into parametric knowledge.

\paragraph{MTP and Context Distillation}
MTP self-distillation \citep{2602.06019} forces multi-step future prediction, regularizing representations with long-horizon self-supervision. Context Distillation \citep{2602.12275} distills the distribution of a heavily prompted model (e.g., ``Think step-by-step'') into the unprompted baseline, baking contextual reasoning into the weights.

\paragraph{Online Experiential Learning}
\citet{2603.16856} maintain a replay buffer of highest-reward rollouts and optimize a dual objective: generating novel trajectories while minimizing KL divergence to successful historical policies, preventing catastrophic forgetting during self-play.

\subsection{Why Self-Play Saturates and Breakthrough Strategies}
\label{subsec:saturation_analysis}

Self-play rapidly saturates after a few iterations due to mode collapse: the model optimizes against a target (itself) sharing the same biases and limitations. If it discovers a confident but flawed heuristic, no external signal corrects it; the model self-reinforces the error until $p_S(y_{flawed}|x) \to 1$, gradients vanish, and exploration ceases.

\textbf{Strategies to Break Through Saturation.} Two approaches inject external truth into the closed loop:
1. \textbf{External Verifiers.} As in OPSD, non-differentiable verifiers (code execution, symbolic math) assign $R=0$ to confident but wrong outputs, breaking the echo chamber.
2. \textbf{The Distillation-RL Loop.} Interleaving self-distillation with PPO-based RL using an independently trained reward model:
\begin{equation}
    \max_\theta \mathbb{E}_{y \sim p_\theta} \left[ R(y) - \beta D_{KL}(p_\theta \parallel p_{ref}) \right]
\end{equation}
Periodically refreshing the reward model on new student generations keeps the target distribution moving, preventing premature convergence.


\section{Reasoning Distillation}
\label{sec:reasoning}

Complex reasoning represents the most active frontier in LLM research. OPD is not merely a compression technique here---it is the mechanism for structured knowledge transfer, shifting training from the teacher's static dataset to the student's active exploration space.

\subsection{Chain-of-Thought Distillation}
As established by \citet{2305.02301}, allowing smaller models to learn from on-policy rationales enables them to outperform un-fine-tuned large models. When the student generates its own reasoning steps, the teacher shifts from dictating outputs to evaluating the student's intrinsic logic. Let $r = (r_1, \dots, r_T)$ be intermediate reasoning tokens. The on-policy CoT objective is:
\begin{align}
\mathcal{L}_{\text{CoT-OPD}}(\theta) &= \mathbb{E}_{x \sim \mathcal{D}, r^S \sim P_\theta(\cdot|x)} \left[ \sum_{t=1}^{|r^S|} \mathcal{D}_{\text{KL}} \Big( P_{\text{Teacher}}(\cdot | x, r^S_{<t}) \parallel P_\theta(\cdot | x, r^S_{<t}) \Big) \right] \label{eq:cot_opd}
\end{align}
Unlike off-policy Forward KL (mean-seeking, which forces the student to cover all teacher modes), the Reverse KL in Equation~\ref{eq:cot_opd} is mode-seeking, reinforcing reasoning paths highly probable under the student's own parameterization when validated by the teacher.

SuperCorrect \citep{2410.09008} augments this with self-correction: the generation space is partitioned into thought $z \sim P_\theta(\cdot|x)$ and answer $y \sim P_\theta(\cdot|x, z)$. If the teacher detects an error in $y$, the student generates a correction $c \sim P_\theta(\cdot|x, z, y)$:
\begin{align}
\mathcal{L}_{\text{SuperCorrect}}(\theta) &= \mathbb{E}_{x, z, y} \left[ \lambda_1 \mathcal{L}_{\text{KD}}(P_T \parallel P_\theta | x, z, y, c) + \lambda_2 \mathcal{R}(c, y_{\text{true}}) \right]
\end{align}
where $\mathcal{R}$ indicates whether the correction recovered the true answer. This teaches the student error detection and recovery, going beyond standard behavioral cloning.

\subsection{Reward-Guided On-Policy Distillation}
RL integration bridges imitation learning and outcome-driven optimization, fusing outcome-based (sparse, flexible) and logit-based (dense, white-box) signals. RLKD \citep{2505.16142} and AlignDistil \citep{2503.02832} unify alignment and distillation by using the teacher's implicit reward to guide on-policy generation:
\begin{align}
\max_{\theta} \mathcal{J}(\theta) = \mathbb{E}_{x \sim \mathcal{D}, y \sim P_\theta(\cdot|x)} \left[ R_T(x, y) - \beta \mathcal{D}_{\text{KL}}\big(P_\theta(\cdot|x) \parallel P_{\text{Teacher}}(\cdot|x)\big) \right] \label{eq:rlkd}
\end{align}
In reward-aware OPD, sparse outcome rewards weight the dense token-level distillation loss via the REINFORCE estimator with a distillation baseline:
\begin{align}
\nabla_\theta \mathcal{J}(\theta) \approx \mathbb{E}_{y \sim P_\theta} \left[ \sum_{t=1}^{|y|} \nabla_\theta \log P_\theta(y_t | x, y_{<t}) \Big( \hat{A}_t(x, y) - \lambda \nabla_\theta \mathcal{D}_{\text{KL}}(P_T \parallel P_\theta) \Big) \right]
\end{align}
where $\hat{A}_t$ is the advantage from the teacher's value function. This enables Reward Extrapolation \citep{2602.12125}: the student explores novel reasoning paths where $P_T(y|x)$ is low but $R_T(x, y)$ is high, escaping the performance ceiling of strict behavioral cloning.

\subsection{The DeepSeek-R1 Paradigm Shift}
DeepSeek-R1 \citep{2501.12948} showed that massive RL on a large model produces highly structured, verifiable reasoning traces with emergent self-verification (the ``Aha moment''). However, applying RL directly to small models ($<$10B parameters) often causes mode collapse or reward hacking due to high gradient variance. The variance comparison is stark:
\begin{align}
\text{Var}(\hat{g}_{\text{RL}}) &\propto \mathbb{E}\left[ \Big( \sum_t \nabla_\theta \log P_\theta(a_t|s_t) A(s_t, a_t) \Big)^2 \right] \\
\text{Var}(\hat{g}_{\text{KD}}) &\propto \mathbb{E}\left[ \Big( \nabla_\theta \mathcal{D}_{\text{KL}}(P_T \parallel P_\theta) \Big)^2 \right] \ll \text{Var}(\hat{g}_{\text{RL}})
\end{align}
Small models lack the representational capacity to act as stable function approximators under high $\text{Var}(\hat{g}_{\text{RL}})$.

DeepSeek-R1 solves this by decoupling discovery from internalization: the large model performs RL-driven exploration, then its reasoning trajectories (and dense token-level distributions) are distilled into smaller models via OPD. This proves that OPD is the optimal variance-reduction vehicle for transferring emergent RL structures: the small model inherits complex planning capabilities while optimizing over a smooth KL landscape rather than a rugged reward surface.

\section{Industrial Systems and Scaling}
\label{sec:systems}

Deploying OPD at industrial scale introduces engineering challenges in compute efficiency, scaling, and architectural compatibility.

\subsection{Large-Scale Deployment}
The Qwen3 technical report \citep{2505.09388} describes industrial-grade OPD using a dynamic ensemble of $K$ teacher models with target distribution:
\begin{align}
\tilde{P}_{\text{ensemble}}(y_t | x, y_{<t}) = \sum_{k=1}^K w_k(x) P_{T_k}(y_t | x, y_{<t}), \quad \text{where } \sum w_k(x) = 1
\end{align}
This smoothens the target and prevents overfitting to a single teacher's idiosyncrasies. Gemma 2 \citep{2408.00118} embeds online KD directly into continuous pre-training to bridge performance gaps between parameter tiers.

For edge devices, MobileLLM \citep{2509.24945} shows that sub-billion parameter models can substantially outperform their class through OPD. MiMo-V2 \citep{2601.02780} combines dense multi-teacher logits with token-level rewards from outcome verifiers, proving that heterogeneous feedback signals are synergistic at scale.

\subsection{Efficiency Innovations}
The primary bottleneck of OPD is the teacher's forward pass on dynamically generated student tokens. A naive implementation requires $O(N)$ autoregressive teacher passes for a student sequence of length $N$.

Speculative KD \citep{2410.11325, 2510.24021} uses the student's generations as speculative drafts. The student generates $K$ tokens; the teacher verifies them in parallel. The speedup depends on the acceptance rate $\alpha$:
\begin{align}
E[S] = \frac{\mathbb{E}[K_{\text{accepted}}] + 1}{1 + c} \quad \text{where } \mathbb{E}[K_{\text{accepted}}] = \sum_{k=1}^K \prod_{i=1}^k \min \left(1, \frac{P_T(y_{t+i} | y_{<t+i})}{P_S(y_{t+i} | y_{<t+i})} \right)
\end{align}
where $c$ is the ratio of student-to-teacher forward pass cost. This accelerates the on-policy loop by 3--5$\times$.

Cross-tokenizer distillation \citep{2402.12030} enables OPD between disparate model families by matching latent semantics. For student vocabulary $V_S$ and teacher vocabulary $V_T$, a projection matrix or optimal transport plan minimizes the Wasserstein distance:
\begin{align}
\mathcal{L}_{\text{CrossTok}} = \inf_{\pi \in \Pi(P_S, P_T)} \mathbb{E}_{(z_S, z_T) \sim \pi} \left[ \parallel W_{S \to T} z_S - z_T \parallel_2^2 \right]
\end{align}
Minitron \citep{2407.14679} combines structured pruning with OPD: a binary mask over teacher weights is updated jointly with the OPD loss, so the student physically emerges from the teacher's architecture, reducing convergence time.

\subsection{Compute-Quality Tradeoff Analysis}
A neglected area is the cost-benefit analysis of OPD. The compute costs (in FLOPs) for $N$ tokens are:
\begin{align}
C_{\text{off}} &\approx N \times (F_{\text{teacher}} + F_{\text{student}} + B_{\text{student}}) \\
C_{\text{on}} &\approx N \times \left( G_{\text{student}} + \lambda F_{\text{teacher}} + F_{\text{student}} + B_{\text{student}} \right)
\end{align}
where $F, B$ denote forward and backward pass FLOPs, $G$ is auto-regressive generation cost (scaling quadratically with sequence length), and $\lambda \in (0, 1]$ is the teacher supervision refresh rate. Since $G_{\text{student}} \gg F_{\text{student}}$, $C_{\text{on}}$ is typically 4--10$\times$ higher than $C_{\text{off}}$.

The constrained optimization is:
\begin{align}
\max \mathcal{U}(\theta, C) \quad \text{subject to } C = \gamma C_{\text{on}} + (1 - \gamma) C_{\text{off}} \le C_{\text{max}}
\end{align}
where $\gamma$ is the on-policy data ratio. For generic knowledge transfer, $\gamma \to 0$ (off-policy) suffices. For reasoning, code, and mathematics, the performance gradient $\frac{\partial \mathcal{U}}{\partial C_{\text{on}}}$ vastly exceeds $\frac{\partial \mathcal{U}}{\partial C_{\text{off}}}$, justifying the compute multiplier. A practical hybrid curriculum starts with off-policy warmup and shifts to $\gamma \to 1$ during the final 20\% of training.

\section{Open Problems and Future Directions}
\label{sec:future}

Despite the empirical success of OPD, theoretical foundations remain sparse. We outline the most critical open problems.

\subsection{Distillation Scaling Laws}
No equivalent of Chinchilla scaling laws exists for OPD. Practitioners guess the optimal generation budget. Understanding how distillation loss scales with $N_S$ (student params), $N_T$ (teacher params), and $D_{\text{on}}$ (on-policy tokens) is vital. A candidate joint scaling law:
\begin{align}
L(N_S, N_T, D_{\text{on}}) = E + \frac{A}{N_S^\alpha} + \frac{B}{N_T^\beta} + \frac{C}{D_{\text{on}}^\gamma} + f(N_S, N_T) \label{eq:scaling}
\end{align}
where $E$ is the irreducible task entropy and $f(N_S, N_T)$ models the capacity gap. Empirically fitting $\alpha, \beta, \gamma$ would let practitioners answer: ``Given $10^4$ GPU hours, distill from 70B for 1B tokens or 405B for 200M tokens?''

\subsection{Uncertainty-Aware On-Policy Distillation}
Teachers hallucinate and exhibit overconfidence in OOD regions. Standard OPD treats the teacher as ground truth, creating an echo chamber. Future solutions should decompose teacher uncertainty into aleatoric and epistemic components. With Bayesian predictive variance $U_{\text{epi}}(x) = \text{Var}_{\phi \sim q}[P_\phi(y|x)]$:
\begin{align}
\mathcal{L}_{\text{Uncertainty-OPD}} = \mathbb{E}_{r \sim P_\theta} \left[ \sum_t \exp\big(-U_{\text{epi}}(r_{<t})\big) \mathcal{D}_{\text{KL}}\Big(P_T(\cdot | r_{<t}) \parallel P_\theta(\cdot | r_{<t})\Big) \right]
\end{align}
Attenuating the gradient when $U_{\text{epi}}$ is high lets the student rely on its own priors rather than mimicking noise.

\subsection{Dynamic Curriculum Distillation}
Uniform prompt sampling wastes compute: easy prompts yield near-zero gradients; hard prompts produce noise. Inspired by PACED \citep{2603.11178}, prompt sampling should evolve with student capability:
\begin{align}
P_t(x) \propto \exp\left( -\frac{| \mathcal{D}_{\text{KL}}(P_T || P_{\theta_t}) - \delta_t |}{\tau} \right)
\end{align}
where $\delta_t$ is an increasing pacing function. This keeps training at the boundary of student capability (the zone of proximal development).

\subsection{Latent Space Distillation for Vocabulary Mismatch}
Equation~\ref{eq:cot_opd} assumes shared vocabulary. In practice, cross-family distillation (e.g., Llama to Phi) involves tokenizer mismatch. The solution is latent space alignment. For final hidden states $h_t^S \in \mathbb{R}^{d_S}$ and $h_t^T \in \mathbb{R}^{d_T}$:
\begin{align}
\mathcal{L}_{\text{Latent}} = \min_{W \in \mathbb{R}^{d_T \times d_S}} \mathbb{E}_{x} \left[ \parallel W h^S(x) - h^T(x) \parallel_2^2 + \lambda \mathcal{R}(W) \right]
\end{align}
where $\mathcal{R}(W)$ enforces orthogonality or sparsity. This bypasses the vocabulary bottleneck, enabling cross-architecture OPD.

\subsection{On-Policy Distillation for Autonomous Agents}
Current OPD targets single-turn generation. LLMs deployed as autonomous agents (using APIs, tools, databases) require long-horizon credit assignment that token-level distillation cannot provide. Agentic OPD should be framed as a POMDP. For trajectory $\tau = (o_1, a_1, \dots, o_T, a_T)$ with teacher value function $Q_T$:
\begin{align}
\mathcal{L}_{\text{Agent-OPD}} = \mathbb{E}_{\tau \sim P_\theta} \left[ \sum_{t=1}^T \mathcal{D}_{\text{KL}}\Big( \pi_T(\cdot | h_t) \parallel \pi_\theta(\cdot | h_t) \Big) \cdot Q_T(h_t, a_t) \right]
\end{align}
This enables distilling strategic foresight and tool-use capabilities, not just linguistic syntax.

\subsection{The Distillation-RL Virtuous Loop}
Distillation and RL are typically applied as sequential stages, but static distillation saturates and pure RL diverges. The optimal approach alternates: RL expands the frontier ($\theta_{k+1/2} = \theta_k + \eta \nabla J_{\text{RL}}$), then a distillation projection pulls back toward the teacher manifold:
\begin{align}
\theta_{k+1} &= \arg\min_{\theta} \left( \parallel \theta - \theta_{k+1/2} \parallel_2^2 + \lambda \mathcal{D}_{\text{KL}}(P_T \parallel P_\theta) \right)
\end{align}
Proving convergence bounds for this alternating operator is a key theoretical direction for recursive self-improvement.

\subsection{Evaluation Methodology Beyond Benchmarks}
Standard benchmarks (MMLU, GSM8K) are contaminated and measure static behavioral cloning, not robustness or OOD generalization. Evaluation should move toward dynamic adversarial testing. We define the generalization gap as:
\begin{align}
\Delta_{\text{OOD}} = \max_{\parallel \epsilon \parallel \le \delta} \text{JS}\Big( P_T(\cdot | x + \epsilon) \parallel P_S(\cdot | x + \epsilon) \Big)
\end{align}
Benchmarking suites should compute this across semantically equivalent but syntactically diverse prompts to verify whether the student has learned causal reasoning mechanisms rather than superficial correlations.



\section{Conclusion}
\label{sec:conclusion}
On-Policy Distillation has rewritten the rules of LLM scaling and deployment. By replacing static dataset mimicry with active, trajectory-level exploration guided by teacher distributions, the community has unlocked reasoning capabilities in constrained architectures previously thought impossible. This survey has dissected the algorithmic taxonomy of OPD---from Chain-of-Thought extraction to Reward-Guided formulations and the DeepSeek-R1 pipeline---formalized the compute-quality tradeoffs, examined innovations enabling large-scale deployment, and charted open theoretical challenges in scaling laws, uncertainty, and latent alignment. As AI shifts toward autonomous agents and recursive self-improvement, OPD will serve as the foundation for the next generation of capable, aligned, and efficient language models.

\section*{Author Contributions}
The authors contributed equally to the formulation, taxonomy design, and analysis of future directions in this survey.

\section*{Acknowledgments}
We acknowledge the open-source LLM community and the authors of the papers surveyed herein.

\bibliography{references}
\bibliographystyle{colm2026_conference}
\end{document}
